% Summary: pending experiments roadmap
% % Summary: pending experiments roadmap
% % Summary: pending experiments roadmap
% % Summary: pending experiments roadmap
% \input{40-summary}

\begin{frame}{待做实验总览}
\small
\textbf{当前报告的定位：}A--I 是现有 Qwen layer-20 样例证据；下一阶段要把这些样例变成原论文关心的可量化评估。

\vspace{0.4em}
\begin{tabularx}{\linewidth}{l>{\raggedright\arraybackslash}p{4.0cm}>{\raggedright\arraybackslash}X}
\toprule
路线 & 依赖 & 目标 \\
\midrule
无训练线 \code{nt_*} & released AV/AR、base model、grader、人审、AR rescoring & 判断 final NLA explanation 是否在 known-truth、evaluation awareness、claim support、toy auditing 中有用 \\
训练相关线 \code{tr_*} & NLA checkpoints、SAE、direct oracle、trained hidden-objective target & 判断能力是否随训练形成，是否支撑原文 automated auditing 强主张 \\
\bottomrule
\end{tabularx}

\vspace{0.6em}
\alert{报告口径：}在这些待做实验完成前，当前 case 只能作为假设生成和实验设计依据，不作为隐藏意图或 claim truth 的最终证明。
\end{frame}


\begin{frame}{无训练优先级}
\scriptsize
\begin{tabularx}{\linewidth}{l>{\raggedright\arraybackslash}p{2.9cm}>{\raggedright\arraybackslash}p{4.6cm}>{\raggedright\arraybackslash}X}
\toprule
实验 & 承接 Case & 核心量化方法 & 完成后能说什么 \\
\midrule
\code{nt_eval_awareness} & A / F & awareness grader；prompt echo control；realism edit；按 sampling window 分层 & evaluation/test hit 是否真的高于格式和 prompt 复述 \\
\code{nt_confabulation_claim} & E & claim truth/support/recurrence 人审；删除单条 claim 后重跑 AR cos/MSE & 哪些 explanation claims 对 reconstruction 有贡献，哪些只是貌似具体 \\
\code{nt_known_truth_eval} & C / D / G / H & 已知标签恢复；候选答案选择；shuffled activation 与 random explanation baseline & NLA 是否携带可由 grader 恢复的任务/语言/答案信息 \\
\code{nt_steganography_quality} & D / E & paraphrase、compression、summary 变换后 AR score drop；writing quality / parse failure & 保真是否依赖非人类可读暗码或过度模板化文本 \\
\code{nt_toy_auditing} & B / I & toy hidden objective wrapper；agent with/without NLA；root-cause grading & NLA 工具链是否能帮助审计 toy hidden constraints \\
\bottomrule
\end{tabularx}
\end{frame}


\begin{frame}{训练相关优先级}
\tiny
\begin{tabularx}{\linewidth}{l>{\raggedright\arraybackslash}p{3.0cm}>{\raggedright\arraybackslash}p{4.7cm}>{\raggedright\arraybackslash}X}
\toprule
实验 & 承接 Case & 核心量化方法 & 完成后能说什么 \\
\midrule
\code{tr_training_curves} & A / G / H & 多 checkpoint；FVE、known-truth、quality 同步记录 & 能力是否随训练提升，还是只学到格式模板 \\
\code{tr_sae_agreement} & E / G & NLA claims 与 SAE top features 对齐 & NLA 主题/具体 claim 是否有外部解释工具支持 \\
\code{tr_direct_activation_oracle} & D & 训练 activation-to-answer oracle；比较 NLA + grader / shuffled & 直接读 activation 是否优于自然语言解释管线 \\
\code{tr_hidden_objective_auditing} & B / I & trained hidden-objective target；win rate 与 root-cause accuracy & 是否支撑原文 automated auditing 应用主张 \\
\bottomrule
\end{tabularx}

\vspace{0.45em}
\alert{优先顺序：}先完成 \code{nt_eval_awareness} 与 \code{nt_confabulation_claim}，再决定 oracle / SAE / target organism 投入。
\end{frame}


\begin{frame}{待做实验总览}
\small
\textbf{当前报告的定位：}A--I 是现有 Qwen layer-20 样例证据；下一阶段要把这些样例变成原论文关心的可量化评估。

\vspace{0.4em}
\begin{tabularx}{\linewidth}{l>{\raggedright\arraybackslash}p{4.0cm}>{\raggedright\arraybackslash}X}
\toprule
路线 & 依赖 & 目标 \\
\midrule
无训练线 \code{nt_*} & released AV/AR、base model、grader、人审、AR rescoring & 判断 final NLA explanation 是否在 known-truth、evaluation awareness、claim support、toy auditing 中有用 \\
训练相关线 \code{tr_*} & NLA checkpoints、SAE、direct oracle、trained hidden-objective target & 判断能力是否随训练形成，是否支撑原文 automated auditing 强主张 \\
\bottomrule
\end{tabularx}

\vspace{0.6em}
\alert{报告口径：}在这些待做实验完成前，当前 case 只能作为假设生成和实验设计依据，不作为隐藏意图或 claim truth 的最终证明。
\end{frame}


\begin{frame}{无训练优先级}
\scriptsize
\begin{tabularx}{\linewidth}{l>{\raggedright\arraybackslash}p{2.9cm}>{\raggedright\arraybackslash}p{4.6cm}>{\raggedright\arraybackslash}X}
\toprule
实验 & 承接 Case & 核心量化方法 & 完成后能说什么 \\
\midrule
\code{nt_eval_awareness} & A / F & awareness grader；prompt echo control；realism edit；按 sampling window 分层 & evaluation/test hit 是否真的高于格式和 prompt 复述 \\
\code{nt_confabulation_claim} & E & claim truth/support/recurrence 人审；删除单条 claim 后重跑 AR cos/MSE & 哪些 explanation claims 对 reconstruction 有贡献，哪些只是貌似具体 \\
\code{nt_known_truth_eval} & C / D / G / H & 已知标签恢复；候选答案选择；shuffled activation 与 random explanation baseline & NLA 是否携带可由 grader 恢复的任务/语言/答案信息 \\
\code{nt_steganography_quality} & D / E & paraphrase、compression、summary 变换后 AR score drop；writing quality / parse failure & 保真是否依赖非人类可读暗码或过度模板化文本 \\
\code{nt_toy_auditing} & B / I & toy hidden objective wrapper；agent with/without NLA；root-cause grading & NLA 工具链是否能帮助审计 toy hidden constraints \\
\bottomrule
\end{tabularx}
\end{frame}


\begin{frame}{训练相关优先级}
\tiny
\begin{tabularx}{\linewidth}{l>{\raggedright\arraybackslash}p{3.0cm}>{\raggedright\arraybackslash}p{4.7cm}>{\raggedright\arraybackslash}X}
\toprule
实验 & 承接 Case & 核心量化方法 & 完成后能说什么 \\
\midrule
\code{tr_training_curves} & A / G / H & 多 checkpoint；FVE、known-truth、quality 同步记录 & 能力是否随训练提升，还是只学到格式模板 \\
\code{tr_sae_agreement} & E / G & NLA claims 与 SAE top features 对齐 & NLA 主题/具体 claim 是否有外部解释工具支持 \\
\code{tr_direct_activation_oracle} & D & 训练 activation-to-answer oracle；比较 NLA + grader / shuffled & 直接读 activation 是否优于自然语言解释管线 \\
\code{tr_hidden_objective_auditing} & B / I & trained hidden-objective target；win rate 与 root-cause accuracy & 是否支撑原文 automated auditing 应用主张 \\
\bottomrule
\end{tabularx}

\vspace{0.45em}
\alert{优先顺序：}先完成 \code{nt_eval_awareness} 与 \code{nt_confabulation_claim}，再决定 oracle / SAE / target organism 投入。
\end{frame}


\begin{frame}{待做实验总览}
\small
\textbf{当前报告的定位：}A--I 是现有 Qwen layer-20 样例证据；下一阶段要把这些样例变成原论文关心的可量化评估。

\vspace{0.4em}
\begin{tabularx}{\linewidth}{l>{\raggedright\arraybackslash}p{4.0cm}>{\raggedright\arraybackslash}X}
\toprule
路线 & 依赖 & 目标 \\
\midrule
无训练线 \code{nt_*} & released AV/AR、base model、grader、人审、AR rescoring & 判断 final NLA explanation 是否在 known-truth、evaluation awareness、claim support、toy auditing 中有用 \\
训练相关线 \code{tr_*} & NLA checkpoints、SAE、direct oracle、trained hidden-objective target & 判断能力是否随训练形成，是否支撑原文 automated auditing 强主张 \\
\bottomrule
\end{tabularx}

\vspace{0.6em}
\alert{报告口径：}在这些待做实验完成前，当前 case 只能作为假设生成和实验设计依据，不作为隐藏意图或 claim truth 的最终证明。
\end{frame}


\begin{frame}{无训练优先级}
\scriptsize
\begin{tabularx}{\linewidth}{l>{\raggedright\arraybackslash}p{2.9cm}>{\raggedright\arraybackslash}p{4.6cm}>{\raggedright\arraybackslash}X}
\toprule
实验 & 承接 Case & 核心量化方法 & 完成后能说什么 \\
\midrule
\code{nt_eval_awareness} & A / F & awareness grader；prompt echo control；realism edit；按 sampling window 分层 & evaluation/test hit 是否真的高于格式和 prompt 复述 \\
\code{nt_confabulation_claim} & E & claim truth/support/recurrence 人审；删除单条 claim 后重跑 AR cos/MSE & 哪些 explanation claims 对 reconstruction 有贡献，哪些只是貌似具体 \\
\code{nt_known_truth_eval} & C / D / G / H & 已知标签恢复；候选答案选择；shuffled activation 与 random explanation baseline & NLA 是否携带可由 grader 恢复的任务/语言/答案信息 \\
\code{nt_steganography_quality} & D / E & paraphrase、compression、summary 变换后 AR score drop；writing quality / parse failure & 保真是否依赖非人类可读暗码或过度模板化文本 \\
\code{nt_toy_auditing} & B / I & toy hidden objective wrapper；agent with/without NLA；root-cause grading & NLA 工具链是否能帮助审计 toy hidden constraints \\
\bottomrule
\end{tabularx}
\end{frame}


\begin{frame}{训练相关优先级}
\tiny
\begin{tabularx}{\linewidth}{l>{\raggedright\arraybackslash}p{3.0cm}>{\raggedright\arraybackslash}p{4.7cm}>{\raggedright\arraybackslash}X}
\toprule
实验 & 承接 Case & 核心量化方法 & 完成后能说什么 \\
\midrule
\code{tr_training_curves} & A / G / H & 多 checkpoint；FVE、known-truth、quality 同步记录 & 能力是否随训练提升，还是只学到格式模板 \\
\code{tr_sae_agreement} & E / G & NLA claims 与 SAE top features 对齐 & NLA 主题/具体 claim 是否有外部解释工具支持 \\
\code{tr_direct_activation_oracle} & D & 训练 activation-to-answer oracle；比较 NLA + grader / shuffled & 直接读 activation 是否优于自然语言解释管线 \\
\code{tr_hidden_objective_auditing} & B / I & trained hidden-objective target；win rate 与 root-cause accuracy & 是否支撑原文 automated auditing 应用主张 \\
\bottomrule
\end{tabularx}

\vspace{0.45em}
\alert{优先顺序：}先完成 \code{nt_eval_awareness} 与 \code{nt_confabulation_claim}，再决定 oracle / SAE / target organism 投入。
\end{frame}


\begin{frame}{待做实验总览}
\small
\textbf{当前报告的定位：}A--I 是现有 Qwen layer-20 样例证据；下一阶段要把这些样例变成原论文关心的可量化评估。

\vspace{0.4em}
\begin{tabularx}{\linewidth}{l>{\raggedright\arraybackslash}p{4.0cm}>{\raggedright\arraybackslash}X}
\toprule
路线 & 依赖 & 目标 \\
\midrule
无训练线 \code{nt_*} & released AV/AR、base model、grader、人审、AR rescoring & 判断 final NLA explanation 是否在 known-truth、evaluation awareness、claim support、toy auditing 中有用 \\
训练相关线 \code{tr_*} & NLA checkpoints、SAE、direct oracle、trained hidden-objective target & 判断能力是否随训练形成，是否支撑原文 automated auditing 强主张 \\
\bottomrule
\end{tabularx}

\vspace{0.6em}
\alert{报告口径：}在这些待做实验完成前，当前 case 只能作为假设生成和实验设计依据，不作为隐藏意图或 claim truth 的最终证明。
\end{frame}


\begin{frame}{无训练优先级}
\scriptsize
\begin{tabularx}{\linewidth}{l>{\raggedright\arraybackslash}p{2.9cm}>{\raggedright\arraybackslash}p{4.6cm}>{\raggedright\arraybackslash}X}
\toprule
实验 & 承接 Case & 核心量化方法 & 完成后能说什么 \\
\midrule
\code{nt_eval_awareness} & A / F & awareness grader；prompt echo control；realism edit；按 sampling window 分层 & evaluation/test hit 是否真的高于格式和 prompt 复述 \\
\code{nt_confabulation_claim} & E & claim truth/support/recurrence 人审；删除单条 claim 后重跑 AR cos/MSE & 哪些 explanation claims 对 reconstruction 有贡献，哪些只是貌似具体 \\
\code{nt_known_truth_eval} & C / D / G / H & 已知标签恢复；候选答案选择；shuffled activation 与 random explanation baseline & NLA 是否携带可由 grader 恢复的任务/语言/答案信息 \\
\code{nt_steganography_quality} & D / E & paraphrase、compression、summary 变换后 AR score drop；writing quality / parse failure & 保真是否依赖非人类可读暗码或过度模板化文本 \\
\code{nt_toy_auditing} & B / I & toy hidden objective wrapper；agent with/without NLA；root-cause grading & NLA 工具链是否能帮助审计 toy hidden constraints \\
\bottomrule
\end{tabularx}
\end{frame}


\begin{frame}{训练相关优先级}
\tiny
\begin{tabularx}{\linewidth}{l>{\raggedright\arraybackslash}p{3.0cm}>{\raggedright\arraybackslash}p{4.7cm}>{\raggedright\arraybackslash}X}
\toprule
实验 & 承接 Case & 核心量化方法 & 完成后能说什么 \\
\midrule
\code{tr_training_curves} & A / G / H & 多 checkpoint；FVE、known-truth、quality 同步记录 & 能力是否随训练提升，还是只学到格式模板 \\
\code{tr_sae_agreement} & E / G & NLA claims 与 SAE top features 对齐 & NLA 主题/具体 claim 是否有外部解释工具支持 \\
\code{tr_direct_activation_oracle} & D & 训练 activation-to-answer oracle；比较 NLA + grader / shuffled & 直接读 activation 是否优于自然语言解释管线 \\
\code{tr_hidden_objective_auditing} & B / I & trained hidden-objective target；win rate 与 root-cause accuracy & 是否支撑原文 automated auditing 应用主张 \\
\bottomrule
\end{tabularx}

\vspace{0.45em}
\alert{优先顺序：}先完成 \code{nt_eval_awareness} 与 \code{nt_confabulation_claim}，再决定 oracle / SAE / target organism 投入。
\end{frame}
